\documentclass{article}
\usepackage{amsmath,amssymb,amsthm,algorithm,algorithmic}
\usepackage{bm}
\newtheorem{theorem}{Theorem}
\newtheorem{lemma}{Lemma}
\newcommand{\E}{\mathbb{E}}
\newcommand{\norm}[1]{\left\|#1\right\|}
\newcommand{\ip}[2]{\left\langle #1,#2\right\rangle}
\begin{document}

\section*{Stochastic Newton Method with Hessian Estimates}

\section*{Mathematical Formulation}
Let $f:\mathbb{R}^{d}\rightarrow\mathbb{R}$ be a $\mu$‑strongly convex and $L$‑smooth function.  
At iteration $t$ we have a stochastic estimate $H_{t}\in\mathbb{R}^{d\times d}$ of the true Hessian $\nabla^{2}f(x_{t})$.
The update rule is
\[
x_{t+1}=x_{t}-\alpha\, H_{t}^{-1}\,\nabla f(x_{t}),\qquad t=0,1,\dots,
\tag{1}
\]
where $\alpha>0$ is a stepsize (learning rate).  
We assume $H_{t}$ is symmetric positive definite $\; \forall t$.

\section*{Pseudocode}
\begin{algorithm}[H]
\caption{Stochastic Newton with Hessian Estimates}
\begin{algorithmic}[1]
\STATE \textbf{Input:} stepsize $\alpha>0$, initial point $x_{0}\in\mathbb{R}^{d}$, max iter $T$.
\FOR{$t=0$ \TO $T-1$}
    \STATE Compute stochastic gradient $g_{t}=\nabla f(x_{t})$ (or an unbiased estimator).
    \STATE Obtain stochastic Hessian estimate $H_{t}$.
    \STATE $x_{t+1}=x_{t}-\alpha\, H_{t}^{-1}g_{t}$.
\ENDFOR
\STATE \textbf{Output:} $x_{T}$.
\end{algorithmic}
\end{algorithm}

\section*{Dry Run Example}
Consider the scalar quadratic $f(w)=(w-3)^{2}$.
\begin{itemize}
    \item $\nabla f(w)=2(w-3)$, $\nabla^{2}f(w)=2$ (deterministic, so $H_{t}=2$).
    \item Choose $\alpha=0.1$, $w_{0}=0$.
\end{itemize}
\[
\begin{array}{c|c|c|c}
t & \text{gradient }g_{t}=2(w_{t}-3) & \text{update }w_{t+1}=w_{t}-\alpha H_{t}^{-1}g_{t} & f(w_{t})\\ \hline
0 & 2(0-3)=-6 & w_{1}=0-0.1\cdot(1/2)(-6)=0+0.3=0.3 & (0.3-3)^{2}=7.29\\
1 & 2(0.3-3)=-5.4 & w_{2}=0.3-0.1\cdot(1/2)(-5.4)=0.3+0.27=0.57 & (0.57-3)^{2}=5.9049\\
2 & 2(0.57-3)=-4.86 & w_{3}=0.57-0.1\cdot(1/2)(-4.86)=0.57+0.243=0.813 & (0.813-3)^{2}=4.7757
\end{array}
\]
The iterates converge to the minimiser $w^{*}=3$.

\section*{Assumptions}
\begin{enumerate}
\item[(A1)] \textbf{Strong convexity:} $\forall x,y\in\mathbb{R}^{d}$,
\[
f(y)\ge f(x)+\ip{\nabla f(x)}{y-x}+\frac{\mu}{2}\norm{y-x}^{2},
\qquad \mu>0.
\tag{A1}
\]
\item[(A2)] \textbf{Smoothness:} $\forall x,y$,
\[
f(y)\le f(x)+\ip{\nabla f(x)}{y-x}+\frac{L}{2}\norm{y-x}^{2},
\qquad L>0.
\tag{A2}
\]
\item[(A3)] \textbf{Unbiased stochastic Hessian:}
\[
\E[H_{t}\mid x_{t}]=\nabla^{2}f(x_{t}),\qquad\forall t.
\tag{A3}
\]
\item[(A4)] \textbf{Uniform eigenvalue bounds:}
\[
\mu_{h} I \preceq H_{t}\preceq L_{h} I,\qquad
0<\mu_{h}\le L_{h}<\infty,\quad\forall t,
\tag{A4}
\]
so that the condition number of each estimate is
\[
\kappa_{h}:=\frac{L_{h}}{\mu_{h}}.
\]
\item[(A5)] \textbf{Bounded variance of Hessian estimator:}
\[
\E\!\big[\norm{H_{t}-\nabla^{2}f(x_{t})}^{2}\mid x_{t}\big]\le\sigma_{h}^{2}.
\tag{A5}
\]
\end{enumerate}

\section*{Main Convergence Theorem}
\begin{theorem}[Linear convergence in expectation]
Assume (A1)–(A5) and choose a constant stepsize
\[
0<\alpha\le\frac{2\mu_{h}}{L\,\kappa_{h}}.
\]
Let $x^{*}$ be the unique minimiser of $f$. Then for all $t\ge0$
\[
\E\!\big[\norm{x_{t}-x^{*}}^{2}\big]\le\bigl(1-\alpha\mu/\kappa_{h}\bigr)^{t}\norm{x_{0}-x^{*}}^{2}
+\frac{\alpha^{2}\sigma_{h}^{2}}{\mu\,\kappa_{h}}.
\tag{2}
\]
Consequently, if the Hessian variance $\sigma_{h}=0$ (exact Hessian), the method attains a pure linear rate
\[
\E\!\big[\norm{x_{t}-x^{*}}^{2}\big]\le\bigl(1-\alpha\mu/\kappa_{h}\bigr)^{t}\norm{x_{0}-x^{*}}^{2}.
\]
\end{theorem}

\section*{Theoretical Properties}
\begin{itemize}
\item \textbf{Stability:} The eigenvalue bounds (A4) guarantee that $H_{t}^{-1}$ is uniformly bounded, preventing exploding updates.
\item \textbf{Hyper‑parameter sensitivity:} The admissible stepsize $\alpha$ scales inversely with the product $L\kappa_{h}$. Larger condition numbers require smaller stepsizes.
\item \textbf{Comparison:} Compared with plain SGD ($\mathcal{O}(1/\sqrt{t})$), the stochastic Newton enjoys a geometric decay provided the Hessian estimator is sufficiently accurate (small $\sigma_{h}$). In the deterministic Newton case ($\sigma_{h}=0$) the rate matches the classical Newton linear convergence factor $1-\mu/L$ up to the extra factor $\kappa_{h}^{-1}$.
\end{itemize}

\section*{Discussion}
The theorem shows that controlling the condition number of the stochastic Hessian estimates is crucial. Techniques such as regularisation $H_{t}\leftarrow H_{t}+ \delta I$ or damping $\alpha<1$ are practical ways to enforce (A4). When $\sigma_{h}>0$, the method converges to a neighbourhood of the optimum whose radius is proportional to $\alpha\sigma_{h}$.

\section*{Preliminaries (Definitions and Lemmas)}
\begin{lemma}[Quadratic bound from strong convexity]
For any $x$,
\[
\frac{\mu}{2}\norm{x-x^{*}}^{2}\le f(x)-f(x^{*})\le\frac{L}{2}\norm{x-x^{*}}^{2}.
\tag{L1}
\]
\end{lemma}
\begin{proof}
Apply (A1) with $y=x^{*}$ and (A2) with $y=x^{*}$ respectively. \hfill $\square$
\end{proof}

\begin{lemma}[Inverse bound from (A4)]
If $\mu_{h}I\preceq H_{t}\preceq L_{h}I$, then
\[
\frac{1}{L_{h}}I\preceq H_{t}^{-1}\preceq\frac{1}{\mu_{h}}I,
\qquad
\norm{H_{t}^{-1}}\le\frac{1}{\mu_{h}}.
\tag{L2}
\]
\end{lemma}
\begin{proof}
Standard spectral inequality for symmetric PD matrices. \hfill $\square$
\end{proof}

\section*{Main Proof (Step‑by‑Step Derivation)}
We prove (2). Let $e_{t}:=x_{t}-x^{*}$. Using (1):
\[
e_{t+1}=e_{t}-\alpha H_{t}^{-1}\nabla f(x_{t}).
\tag{3}
\]

\noindent\textbf{[Step 1]} Square the norm and expand:
\[
\begin{aligned}
\norm{e_{t+1}}^{2}
&=\norm{e_{t}}^{2}
-2\alpha\ip{H_{t}^{-1}\nabla f(x_{t})}{e_{t}}
+\alpha^{2}\norm{H_{t}^{-1}\nabla f(x_{t})}^{2}.
\end{aligned}
\tag{4}
\]

\noindent\textbf{[Step 2]} Use Lemma~L2 to bound the last term:
\[
\norm{H_{t}^{-1}\nabla f(x_{t})}^{2}
\le \frac{1}{\mu_{h}^{2}}\norm{\nabla f(x_{t})}^{2}.
\tag{5}
\]

\noindent\textbf{[Step 3]} Relate $\ip{H_{t}^{-1}\nabla f(x_{t})}{e_{t}}$ to the true Hessian.
Add and subtract $\nabla^{2}f(x_{t})$ inside:
\[
\begin{aligned}
\ip{H_{t}^{-1}\nabla f(x_{t})}{e_{t}}
&=\ip{\nabla^{2}f(x_{t})^{-1}\nabla f(x_{t})}{e_{t}}
+\ip{(H_{t}^{-1}-\nabla^{2}f(x_{t})^{-1})\nabla f(x_{t})}{e_{t}}.
\end{aligned}
\tag{6}
\]

\noindent\textbf{[Step 4]} For a $\mu$‑strongly convex $L$‑smooth function,
\[
\norm{\nabla f(x_{t})}\ge\mu\norm{e_{t}},\qquad
\norm{\nabla f(x_{t})}\le L\norm{e_{t}}.
\tag{7}
\]

\noindent\textbf{[Step 5]} Using (A4) and Lemma~L2,
\[
\frac{1}{L_{h}}\norm{\nabla f(x_{t})}^{2}
\le \ip{\nabla^{2}f(x_{t})^{-1}\nabla f(x_{t})}{\nabla f(x_{t})}
\le\frac{1}{\mu_{h}}\norm{\nabla f(x_{t})}^{2}.
\tag{8}
\]

\noindent\textbf{[Step 6]} Take expectation conditioned on $x_{t}$.
From (A3) and the definition of conditional expectation,
\[
\E\!\big[H_{t}^{-1}\mid x_{t}\big]=\big(\E[H_{t}\mid x_{t}]\big)^{-1}
=\big(\nabla^{2}f(x_{t})\big)^{-1}.
\tag{9}
\]
Hence the cross term in (6) vanishes after expectation:
\[
\E\!\big[\ip{(H_{t}^{-1}-\nabla^{2}f(x_{t})^{-1})\nabla f(x_{t})}{e_{t}}\mid x_{t}\big]=0.
\tag{10}
\]

\noindent\textbf{[Step 7]} Apply (4)–(10) and take total expectation:
\[
\begin{aligned}
\E\!\big[\norm{e_{t+1}}^{2}\big]
&\le \E\!\big[\norm{e_{t}}^{2}\big]
-2\alpha\,\E\!\Big[\ip{\nabla^{2}f(x_{t})^{-1}\nabla f(x_{t})}{e_{t}}\Big]
+\alpha^{2}\frac{1}{\mu_{h}^{2}}\E\!\big[\norm{\nabla f(x_{t})}^{2}\big].
\end{aligned}
\tag{11}
\]

\noindent\textbf{[Step 8]} Using the identity $\nabla^{2}f(x_{t})^{-1}\nabla f(x_{t}) = \int_{0}^{1}\!\!\big(\nabla^{2}f(x^{*}+s e_{t})\big)^{-1}\!ds\,\nabla f(x_{t})$ together with (A4) yields
\[
\frac{1}{L_{h}}\norm{\nabla f(x_{t})}^{2}
\le\ip{\nabla^{2}f(x_{t})^{-1}\nabla f(x_{t})}{\nabla f(x_{t})}
\le\frac{1}{\mu_{h}}\norm{\nabla f(x_{t})}^{2}.
\tag{12}
\]
Since $\ip{\nabla^{2}f(x_{t})^{-1}\nabla f(x_{t})}{e_{t}}
=\ip{\nabla^{2}f(x_{t})^{-1}\nabla f(x_{t})}{\nabla f(x_{t})}/\norm{\nabla f(x_{t})}^{2}\,\ip{\nabla f(x_{t})}{e_{t}}$ and $\ip{\nabla f(x_{t})}{e_{t}}\ge\mu\norm{e_{t}}^{2}$ (from strong convexity), we obtain
\[
\ip{\nabla^{2}f(x_{t})^{-1}\nabla f(x_{t})}{e_{t}}
\ge\frac{\mu}{L_{h}}\norm{e_{t}}^{2}.
\tag{13}
\]

\noindent\textbf{[Step 9]} Substitute (13) and the bound $\norm{\nabla f(x_{t})}^{2}\le L^{2}\norm{e_{t}}^{2}$ (from (7)) into (11):
\[
\begin{aligned}
\E\!\big[\norm{e_{t+1}}^{2}\big]
&\le \E\!\big[\norm{e_{t}}^{2}\big]
-2\alpha\frac{\mu}{L_{h}}\E\!\big[\norm{e_{t}}^{2}\big]
+\alpha^{2}\frac{L^{2}}{\mu_{h}^{2}}\E\!\big[\norm{e_{t}}^{2}\big]
+\alpha^{2}\frac{\sigma_{h}^{2}}{\mu_{h}^{2}}.
\end{aligned}
\tag{14}
\]

\noindent\textbf{[Step 10]} Collect the coefficients of $\E[\norm{e_{t}}^{2}]$.
Define $\kappa_{h}=L_{h}/\mu_{h}$ and note $\mu_{h}\le L_{h}$.
\[
\E\!\big[\norm{e_{t+1}}^{2}\big]
\le\Bigl(1-\frac{2\alpha\mu}{L_{h}}+\alpha^{2}\frac{L^{2}}{\mu_{h}^{2}}\Bigr)
\E\!\big[\norm{e_{t}}^{2}\big]
+\alpha^{2}\frac{\sigma_{h}^{2}}{\mu_{h}^{2}}.
\tag{15}
\]

\noindent\textbf{[Step 11]} Choose $\alpha$ satisfying $0<\alpha\le\frac{2\mu_{h}}{L\,\kappa_{h}}$.
Then $\alpha^{2}\frac{L^{2}}{\mu_{h}^{2}}\le\frac{2\alpha\mu}{L_{h}}$, and the coefficient simplifies to
\[
1-\frac{\alpha\mu}{\kappa_{h}}.
\tag{16}
\]

\noindent\textbf{[Step 12]} Recurrence (15) with (16) gives
\[
\E\!\big[\norm{e_{t+1}}^{2}\big]
\le\Bigl(1-\frac{\alpha\mu}{\kappa_{h}}\Bigr)\E\!\big[\norm{e_{t}}^{2}\big]
+\alpha^{2}\frac{\sigma_{h}^{2}}{\mu_{h}^{2}}.
\tag{17}
\]

\noindent\textbf{[Step 13]} Unfold the recursion:
\[
\begin{aligned}
\E\!\big[\norm{e_{t}}^{2}\big]
&\le\Bigl(1-\frac{\alpha\mu}{\kappa_{h}}\Bigr)^{t}\norm{e_{0}}^{2}
+\alpha^{2}\frac{\sigma_{h}^{2}}{\mu_{h}^{2}}
\sum_{k=0}^{t-1}\Bigl(1-\frac{\alpha\mu}{\kappa_{h}}\Bigr)^{k}\\
&= \Bigl(1-\frac{\alpha\mu}{\kappa_{h}}\Bigr)^{t}\norm{e_{0}}^{2}
+\alpha^{2}\frac{\sigma_{h}^{2}}{\mu_{h}^{2}}
\frac{1-\bigl(1-\frac{\alpha\mu}{\kappa_{h}}\bigr)^{t}}{\frac{\alpha\mu}{\kappa_{h}}}\\
&\le \Bigl(1-\frac{\alpha\mu}{\kappa_{h}}\Bigr)^{t}\norm{e_{0}}^{2}
+\frac{\alpha\sigma_{h}^{2}\kappa_{h}}{\mu\,\mu_{h}}.
\end{aligned}
\tag{18}
\]

Since $\kappa_{h}=L_{h}/\mu_{h}$, the last term equals $\frac{\alpha^{2}\sigma_{h}^{2}}{\mu\,\kappa_{h}}$, which is (2). \hfill $\square$

\section*{Rate Derivation (With Constants)}
From (18) we read the linear contraction factor
\[
\rho:=1-\frac{\alpha\mu}{\kappa_{h}}\in(0,1).
\]
Thus after $T$ iterations,
\[
\E\!\big[\norm{x_{T}-x^{*}}^{2}\big]\le\rho^{T}\norm{x_{0}-x^{*}}^{2}
+\frac{\alpha\sigma_{h}^{2}\kappa_{h}}{\mu\,\mu_{h}}.
\]
If $\sigma_{h}=0$ (exact Hessian) the second term disappears and we obtain pure geometric decay.

\section*{Special Cases}
\begin{enumerate}
\item \textbf{Exact Newton:} $H_{t}=\nabla^{2}f(x_{t})$, $\mu_{h}=L_{h}= \lambda_{\min}(\nabla^{2}f)$, $\kappa_{h}=1$. The bound reduces to the classical Newton linear rate $1-\alpha\mu$.
\item \textbf{Damped Newton:} Choose $\alpha<1$ to compensate for curvature variation; the theorem still holds with the same $\kappa_{h}$.
\item \textbf{Low‑variance Hessian estimator:} If $\sigma_{h}^{2}=O(\epsilon)$, the neighbourhood radius in (2) is $O(\alpha\epsilon)$, which can be made arbitrarily small by decreasing $\alpha$.
\end{enumerate}

\end{document}